\section{Review Questions}

\begin{enumerate}
\item Explain why a matrix should be read as structured information rather than merely as a rectangular table of numbers. Give three different roles that the same array might play.

\item For a matrix $A=(a_{ij})$ of size $m\times n$, explain the meanings of $m$, $n$, $i$, and $j$. Why is the order of the two indices important?

\item Determine whether each operation is defined and state the size of the result:
\begin{enumerate}
\item addition of a $2\times3$ and a $2\times3$ matrix;
\item multiplication of a $2\times3$ by a $3\times4$ matrix;
\item multiplication in the reverse order.
\end{enumerate}
Explain each answer structurally rather than only quoting a rule.

\item Compare entrywise operations with matrix multiplication. What information is preserved in the first case, and what interaction is created in the second?

\item Compute one product of two simple $2\times2$ matrices and explain what each entry of the result records in terms of a row and a column.

\item Why is matrix multiplication generally not commutative? Give a conceptual explanation and verify it with a small example.

\item Explain the roles of the zero matrix and the identity matrix. In what sense do they resemble the numbers $0$ and $1$, and where does the analogy require care?

\item What information can be recognised before calculation in a diagonal, triangular, or symmetric matrix? Give one useful consequence of each visible structure.

\item List the three elementary row operations and explain why each is reversible. Why is reversibility essential when row operations are used to study a problem?

\item A row operation has an arithmetic, representational, and interpretive layer. Explain these three layers using one concrete row replacement and state what relation the new row exposes.
\end{enumerate}
